\documentclass[10pt,a4paper]{amsart}
\usepackage[margin=19mm]{geometry}
\usepackage{amsmath,amssymb,mathtools}
\usepackage[scr=rsfs,cal=euler]{mathalfa}
\usepackage{xfrac}
\usepackage[unicode,hidelinks,bookmarks=false]{hyperref}
\newcommand{\ie}{i.e.\ }
\newcommand{\cf}{cf.\ }
\newcommand{\resp}{respectively\ }
\newcommand{\Subsection}{\subsection}
\providecommand{\nobreakdash}{\nobreak\hbox{-}}
\setlength{\emergencystretch}{2em}
\multlinegap0pt
\usepackage{mathtools}
\usepackage{amssymb}
\usepackage{breqn}
\usepackage[shortlabels]{enumitem}
\usepackage{xcolor}
\newtheorem{theorem}{Theorem}
\newcommand*{\cp}[2]{#1 \,\square\, #2}    
\title{A lattice structure for ancestral configurations arising from the relationship between gene trees and species trees}
\author{Egor Lappo$^*$ and Noah A.\ Rosenberg}
\date{Source: arXiv:2111.10456v3 (CC BY 4.0)}
\begin{document}
\maketitle

\noindent\emph{Source and licence.} A lattice structure for ancestral configurations arising from the relationship between gene trees and species trees. Version arXiv:2111.10456v3 (CC BY 4.0), distributed by arXiv under the Creative Commons Attribution 4.0 International licence, http://creativecommons.org/licenses/by/4.0/, as recorded in the arXiv licence metadata for this exact version and shown on the abstract page https://arxiv.org/abs/2111.10456v3. Full licence notice: \texttt{licenses/arxiv-2111.10456-CC-BY-4.0.txt}.\par\smallskip

\noindent\textit{Editorial scope.} Counted unit: the article's theorem thm:graph\_decomp (source0.tex lines 400--405). The article's complete original proof of it is delivered with it (source0.tex lines 409--417, one \texttt{proof} environment of 9 lines, which the article places away from the statement). No mathematics is written, repaired or re-derived: the statement and the proof are the article's own lines, in the article's own notation, and no retained line is substituted or re-flowed.  Retained uncounted as the source-local context the proof needs: lines 368--375.  Nothing was omitted from any retained span, so the package carries no omission record.  No retained line contains a citation, so the wrapper carries no bibliography and no bibliography entry.  No retained line contains a figure environment, an image-inclusion command or drawing code, so no image is delivered and the package declares no asset dependency.  Editorial formatting only, in addition: an AMS \texttt{amsart} preamble that restates the article's own declarations and macros used by the retained lines, namely the environments \texttt{theorem} on separate counters, together with the standard packages those lines need.  The retained environments print in this document as Theorem 1 in order of appearance, whereas the article prints the same items with the numbering of its own full document.  The complete native source of the article as distributed in the arXiv e-print for this exact version is bundled unchanged as \texttt{sources/118681/source0.tex}.
\begin{equation}
\label{eq:recur_enum}
	C(S) = \{ 
	\left[C(S|_{r_\ell}) \otimes C(S|_{r_r})\right] \cup 
	\left[\{\{r_\ell\}\} \otimes C(S|_{r_r})\right] \cup
    \left[C(S|_{r_\ell}) \otimes \{\{r_r\}\}\right] \cup 
	\{r_\ell,r_r\}\},
\end{equation}

\begin{theorem}\label{thm:graph_decomp}
Consider a tree $S$ with $n\geq 2$ leaves. Let $r_\ell$ and $r_r$ be the two immediate descendants of its root, and let $S|_{r_\ell}$ and $S|_{r_r}$ be the subtrees of $S$ rooted at $r_\ell$ and $r_r$, respectively. Then the digraph $D(S)$ is isomorphic to a cartesian product of $\widehat{D(S|_{r_\ell})}$ and $\widehat{D(S|_{r_r})}$,
\[
	D(S) \cong \cp{\widehat{D(S|_{r_\ell})}}{\widehat{D(S|_{r_r})}}.
\]
\end{theorem}

\begin{proof}[Proof of Theorem~\ref{thm:graph_decomp}]
Eq.~\ref{eq:recur_enum} gives a bijection between the vertex sets $V\big(D(S)\big)$ and $V\left(\cp{\widehat{D(S|_{r_\ell})}}{\widehat{D(S|_{r_r})}}\right)$ by interpreting caterpillar-like maximal nodes added by the ``hat''-operation as $r_r$ and $r_\ell$ in eq.~\ref{eq:recur_enum}. This bijection allows us to formally identify the vertices of the two digraphs.

To conclude that $D(S) = \cp{\widehat{D(S|_{r_\ell})}}{\widehat{D(S|_{r_r})}}$, we must show that the sets of edges of digraphs are equal. We first show containment of edge sets, and then prove by contradiction that containment cannot be strict.

We claim that edge sets satisfy $E\left(\cp{\widehat{D(S|_{r_\ell})}}{\widehat{D(S|_{r_r})}}\right) \subseteq E\big(D(S)\big)$ under the identification of vertices via eq.~\ref{eq:recur_enum}. The digraph $\cp{\widehat{D(S|_{r_\ell})}}{\widehat{D(S|_{r_r})}}$ has an edge between $(u,v)$ and $(u',v')$ if either (1) $u=u'$ and $\widehat{D(S|_{r_r})}$ has an edge from $v$ to $v'$ or (2) if $v = v'$ and $\widehat{D(S|_{r_\ell})}$ has an edge from $u$ to $u'$. This means that each edge of $\cp{\widehat{D(S|_{r_\ell})}}{\widehat{D(S|_{r_r})}}$, viewed as an edge between vertices of $D(S)$, corresponds to precisely one coalescence, happening either in the right subtree $S_r$ (case 1), or in the left subtree $S_\ell$ (case 2), including the root coalescences of the two subtrees. Such an edge is in $E\big(D(S) \big)$ by definition of $D(S)$.

Finally, if there is an edge $e \in E\big(D(S)\big) - E\left(\cp{\widehat{D(S|_{r_\ell})}}{\widehat{D(S|_{r_r})}}\right)$, then it corresponds to a coalescence that happens in neither of the root subtrees $S_\ell$, $S_r$. There is only one such coalescence, namely the root of $S$. However, by definition of $D(S)$, root coalescences are not reflected in the digraph, so we have a contradiction. We conclude that $D(S) = \cp{\widehat{D(S|_{r_\ell})}}{\widehat{D(S|_{r_r})}}$.
\end{proof}

\end{document}
